\chapter{Indivíduos singulares}\label{ch:g9:unique-individuals}

Sete bilhões de pessoas, e nenhuma igual à outra --- nem numa
família, nem num pátio de escola, nem em toda a história, com
os gêmeos idênticos meio excetuados. Três capítulos já
montaram a maquinaria inteira desse fato. Este capítulo curto
de fechamento roda o argumento uma vez, de ponta a ponta, e
guarda as consequências dele: para a variedade da \omterm{def:g5:biodiversity:species}{espécie},
para o modo como lemos as manchetes sobre ``genética'' e para
o que a singularidade significa e não significa.

\section{O teorema da singularidade}

\begin{proposition}[Por que duas pessoas nunca são iguais]\label{prop:g9:unique-individuals:theorem}
Junte os três mecanismos:
\begin{enumerate}
  \item \emph{o embaralhamento da redução à metade}: os
        \omterm{def:g8:sexual-reproduction:gametes}{gametas} de cada um dos pais carregam um parceiro por
        par, distribuído de forma independente --- $2^{23}$
        meios baralhos por cada um deles
        (\cref{rem:g9:chromosomes:whyhalve}) e, na verdade,
        muito mais, já que os parceiros de um par também
        \emph{trocam trechos} durante a divisão de redução à
        metade, cortando o baralho abaixo do nível dos
        \omterm{def:g9:chromosomes:chromosomes}{cromossomos} inteiros;
  \item \emph{a loteria da \omterm{def:g5:flower-to-fruit:steps}{fecundação}}: qual \omterm{def:g8:reproductive-systems:male}{espermatozoide}
        entre os milhões (\cref{ex:g8:reproductive-systems:sperm})
        encontra o óvulo de qual mês multiplica um espaço
        pelo outro;
  \item \emph{o gotejar das \omterm{prop:g9:genes-and-dna:explains}{mutações}}: cada geração
        acrescenta algumas grafias novas dela mesma
        (\cref{prop:g9:genes-and-dna:explains}).
\end{enumerate}
O espaço de combinações faz o número de humanos que já
viveram parecer minúsculo: cada \omterm{def:g5:flower-to-fruit:steps}{fecundação} distribui um
baralho que nunca foi distribuído antes, e nunca mais será.
\end{proposition}
\admitted

\begin{example}[Os gêmeos, e os limites da exceção]\label{ex:g9:unique-individuals:twins}
Os gêmeos idênticos compartilham uma distribuição só
(\cref{exo:g8:from-egg-to-birth:10}) --- e até eles divergem:
vidas separadas escrevem \omterm{def:g9:heredity-and-traits:traits}{características} adquiridas
diferentes (\cref{prop:g9:heredity-and-traits:sorting}),
histórias de \omterm{def:g8:nervous-communication:synapse}{sinapses} diferentes
(\cref{ex:g8:nervous-communication:juggler}) e até algumas
\omterm{prop:g9:genes-and-dna:explains}{mutações} tardias diferentes nas \omterm{def:g5:first-look-at-cells:cell}{células} de cada corpo. A
única rachadura na singularidade se fecha sozinha: depois da
distribuição, o viver diferencia.
\end{example}

\begin{remark}[Singulares, e quase todos iguais]\label{rem:g9:unique-individuals:alike}
Segure as duas verdades: os textos de \omterm{def:g9:genes-and-dna:dna}{DNA} de dois humanos
quaisquer concordam em mais de $99$ por cento --- a
construção compartilhada da \omterm{def:g5:biodiversity:species}{espécie} única da
\cref{prop:g5:first-look-at-cells:principle} --- e as
discordâncias, espalhadas por três bilhões de letras, ainda
bastam para uma singularidade estrita. ``Todos diferentes'' e
``todos parentes'' são o mesmo fato em dois aumentos
(\cref{met:g9:genes-and-dna:levels}); todo uso que já se fez
da primeira verdade contra a segunda era \omterm{def:g1:living-or-not:living}{biologia} ruim antes
de ser ética ruim.
\end{remark}

\section{A singularidade em ação}

\begin{example}[A identificação por DNA]\label{ex:g9:unique-individuals:forensics}
A singularidade dá para conferir, e os tribunais a conferem:
um número suficiente das posições variáveis de uma pessoa,
lidas de um vestígio de \omterm{def:g5:first-look-at-cells:cell}{células}, combina com um indivíduo (à
parte os gêmeos idênticos) em toda a humanidade. A mesma
comparação, rodada no aumento da família, resolve a
paternidade --- cada posição variável do filho tem de poder
ser distribuída a partir dos baralhos dos dois pais. A lógica
do álbum (\cref{ch:g9:heredity-and-traits}) virou um
instrumento.
\end{example}

\begin{example}[A transfusão e o transplante]\label{ex:g9:unique-individuals:medicine}
A medicina negocia com a singularidade todos os dias. Os
grupos sanguíneos (\cref{ex:g9:genes-and-dna:bloodgroups})
são o caso fácil --- quatro classes amplas, combináveis pelo
registro. Os transplantes de órgãos encontram o caso difícil:
o sistema imunológico (\cref{ch:g9:immune-defenses} explica a
maquinaria) lê marcadores de identidade mais finos,
singulares o bastante para que os doadores sejam procurados
primeiro na família --- quanto mais próximo o \omterm{prop:g6:classification-kinship:kinship}{parentesco},
mais próxima a distribuição.
\end{example}

\begin{proposition}[A variedade é a reserva da espécie, guardada]\label{prop:g9:unique-individuals:reserve}
A \Cref{rem:g8:sexual-reproduction:reserve} prometeu o
sentido do embaralhamento; os mecanismos agora o pagam. Uma
\omterm{def:g5:biodiversity:species}{espécie} de indivíduos singulares guarda, espalhados entre
eles, \omterm{def:g9:genes-and-dna:gene}{alelos} para toda ocasião --- a resistência que uns
carregam contra esta doença, a tolerância que outros carregam
àquele estresse (o pomar do
\cref{ex:g8:sexual-reproduction:bets}, na escala da \omterm{def:g5:biodiversity:species}{espécie}).
Quando as \omterm{def:g6:exploring-our-environment:conditions}{condições} viram, os freios da
\cref{prop:g8:reproduction-and-environment:limits} caem de
forma desigual sobre essa variedade --- e o que essa
desigualdade faz, geração após geração, é exatamente o
assunto que os dois próximos capítulos abrem: a mudança das
\omterm{def:g5:biodiversity:species}{espécies}.
\end{proposition}
\admitted

\begin{method}[Auditando uma alegação ``genética'']\label{met:g9:unique-individuals:claims}
Para qualquer manchete sobre \omterm{def:g9:genes-and-dna:gene}{genes} e pessoas:
\begin{enumerate}
  \item conferência de nível
        (\cref{met:g9:genes-and-dna:levels}): a alegação é
        sobre um \omterm{def:g9:genes-and-dna:gene}{gene}, sobre um baralho ou sobre um padrão de
        família?
  \item conferência do tecido
        (\cref{prop:g9:heredity-and-traits:sorting}): ela
        respeita o tecido da faixa herdada preenchida pela
        vida, ou finge que um \omterm{def:g9:genes-and-dna:gene}{gene} é igual a um destino?
  \item conferência da variedade: ela trata a \omterm{def:g6:how-plants-spread:dispersal}{dispersão} de
        uma \omterm{def:g8:reproduction-and-environment:population}{população} como a reserva que ela é, ou como
        desvios de um baralho ``normal'' que não existe?
  \item conferência dos gêmeos: ela sobreviveria aos gêmeos
        idênticos, que compartilham a distribuição e mesmo
        assim diferem?
\end{enumerate}
\end{method}

\section{Exercícios}

\begin{exercise}[$\star$]\label{exo:g9:unique-individuals:1}
Nomeie os três mecanismos do teorema da singularidade.
\end{exercise}

\begin{exercise}[$\star$]\label{exo:g9:unique-individuals:2}
O que a troca de trechos durante a redução à metade
acrescenta ao embaralhamento?
\end{exercise}

\begin{exercise}[$\star$]\label{exo:g9:unique-individuals:3}
Enuncie as duas metades da
\cref{rem:g9:unique-individuals:alike} --- a porcentagem e a
singularidade --- numa frase só.
\end{exercise}

\begin{exercise}[$\star$]\label{exo:g9:unique-individuals:4}
Como os gêmeos idênticos acabam distinguíveis, apesar de tudo?
\end{exercise}

\begin{exercise}[$\star$]\label{exo:g9:unique-individuals:5}
Que duas perguntas a comparação de \omterm{def:g9:genes-and-dna:dna}{DNA} responde para um
tribunal?
\end{exercise}

\begin{exercise}[$\star$]\label{exo:g9:unique-individuals:6}
Por que os doadores de transplante são procurados primeiro na
família?
\end{exercise}

\begin{exercise}[$\star\star$]\label{exo:g9:unique-individuals:7}
Multiplique o argumento: combine os espaços de meios baralhos
dos dois pais e a loteria da \omterm{def:g5:flower-to-fruit:steps}{fecundação} na conclusão do
``nunca distribuído antes''.
\end{exercise}

\begin{exercise}[$\star\star$]\label{exo:g9:unique-individuals:8}
Pague a \cref{prop:g9:unique-individuals:reserve} no vinhedo
do \cref{exo:g8:sexual-reproduction:10}: o que faltava aos
clones que uma \omterm{def:g8:reproduction-and-environment:population}{população} nascida de \omterm{prop:g3:flowers-fruits-seeds:transform}{sementes} tem?
\end{exercise}

\begin{exercise}[$\star\star$]\label{exo:g9:unique-individuals:9}
Faça a \cref{met:g9:unique-individuals:claims} em:
``descoberto um \omterm{def:g9:genes-and-dna:gene}{gene} do sucesso escolar''.
\end{exercise}

\begin{exercise}[$\star\star$]\label{exo:g9:unique-individuals:10}
Por que ``o baralho humano normal'' é uma expressão sem
referente? Responda com a reserva e com o teorema.
\end{exercise}

\begin{exercise}[$\star\star$]\label{exo:g9:unique-individuals:11}
Um teste de paternidade inocenta um pai posto em dúvida e
acusado injustamente pela vizinha do
\cref{exo:g9:heredity-and-traits:11}. Explique o que o teste
confere, posição por posição.
\end{exercise}

\begin{exercise}[$\star\star\star$]\label{exo:g9:unique-individuals:12}
``Todos diferentes, todos parentes, e as duas coisas pelo
mesmo mecanismo.'' Escreva o parágrafo de fechamento:
embaralhamento, \omterm{prop:g9:genes-and-dna:explains}{mutação} e código compartilhado, cada um no
lugar dele.
\end{exercise}

\section{Problema: O seminário de casos arquivados}

\begin{problem}[{Problema de fim de semana --- um fio de
cabelo, três perguntas, trinta anos}]\label{pb:g9:unique-individuals:1}
Um seminário (fictício) sobre casos arquivados: um furto sem
solução de trinta anos atrás, um fio de cabelo conservado com
as \omterm{def:g5:first-look-at-cells:cell}{células} da \omterm{def:g1:plants-around-us:parts}{raiz} dele e os métodos modernos de comparação.
O seminário leva os alunos pelo que a identificação por \omterm{def:g9:genes-and-dna:dna}{DNA}
consegue e não consegue dizer.

\medskip\noindent\textbf{Parte I --- O que o fio de cabelo
guarda.}
\begin{enumerate}
  \item Por que a \emph{\omterm{def:g1:plants-around-us:parts}{raiz}} do cabelo tem de estar
        conservada para uma leitura completa (a
        \cref{def:g6:cell-unit-of-life:parts} sabe onde mora
        o \omterm{def:g9:genes-and-dna:dna}{DNA})?
  \item O laboratório lê um conjunto de posições muito
        variáveis. Por que as variáveis --- o que as posições
        compartilhadas por 99 por cento estabeleceriam?
  \item O perfil combina com um suspeito em todas as
        posições. Diga o que a combinação afirma, e a única
        exceção sistemática dela
        (\cref{ex:g9:unique-individuals:twins}).
  \item A defesa observa que o suspeito tem um irmão gêmeo
        idêntico no exterior. O que o tribunal precisa fazer
        agora, e por quê?
\end{enumerate}

\medskip\noindent\textbf{Parte II --- As perguntas de
família.}
\begin{enumerate}[resume]
  \item Um perfil parcial e anterior tinha apontado para
        ``um parente próximo da família K''. Explique como
        uma combinação \emph{parcial} se lê como \omterm{prop:g6:classification-kinship:kinship}{parentesco}
        (a distribuição da
        \cref{prop:g9:unique-individuals:theorem}, rodada de
        trás para a frente).
  \item Um jurado pergunta se o perfil revela as
        \omterm{def:g9:heredity-and-traits:traits}{características} do suspeito --- aparência,
        temperamento, ``\omterm{def:g9:genes-and-dna:gene}{genes} de criminoso''. Faça a
        \cref{met:g9:unique-individuals:claims} sobre a
        pergunta, conferência por conferência.
  \item Outro jurado pergunta se a combinação não poderia
        ser ``uma vez nesta cidade'', em vez de uma vez na
        humanidade. O que define o número de posições lidas?
  \item A coordenadora do seminário avisa: ``a \omterm{def:g1:living-or-not:living}{biologia} diz
        de quem são as \omterm{def:g5:first-look-at-cells:cell}{células}; o caso ainda precisa dizer
        como elas foram parar ali''. Distinga as duas
        afirmações.
\end{enumerate}

\medskip\noindent\textbf{Parte III --- O encerramento do
seminário.}
\begin{enumerate}[resume]
  \item Resuma o instrumento numa frase só: que mecanismos
        de dois capítulos (\cref{ch:g9:chromosomes},
        \cref{ch:g9:genes-and-dna}) tornam um perfil
        singular e uma família legível.
  \item Por que o mesmo instrumento funciona num fio de
        cabelo de trinta anos atrás? (Que propriedade do \omterm{def:g9:genes-and-dna:dna}{DNA}
        --- à parte a cópia da
        \cref{prop:g9:genes-and-dna:explains} --- está sendo
        usada: a estabilidade.)
  \item O seminário termina na ética: liste dois usos da
        singularidade que este capítulo mostra a serviço das
        pessoas, e o mau uso histórico contra o qual a
        \cref{rem:g9:unique-individuals:alike} adverte.
  \item Encerre com o teorema reenunciado para o tribunal: o
        que nunca aconteceu duas vezes, e por quê.
\end{enumerate}
\end{problem}
